\documentclass{article}

\usepackage{amsmath}
\usepackage{amssymb}
\usepackage{xcolor}
\usepackage{tikz}

\usepackage[utf8]{inputenc}
\usepackage[T1]{fontenc}

% \usepackage[polish]{babel}

\usepackage{listings}
\lstset{
    language=SQL,
    inputencoding=utf8x,
    extendedchars=\true,
    literate={ą}{{\k{a}}}1
             {Ą}{{\k{A}}}1
             {ę}{{\k{e}}}1
             {Ę}{{\k{E}}}1
             {ó}{{\'o}}1
             {Ó}{{\'O}}1
             {ś}{{\'s}}1
             {Ś}{{\'S}}1
             {ł}{{\l{}}}1
             {Ł}{{\L{}}}1
             {ż}{{\.z}}1
             {Ż}{{\.Z}}1
             {ź}{{\'z}}1
             {Ź}{{\'Z}}1
             {ć}{{\'c}}1
             {Ć}{{\'C}}1
             {ń}{{\'n}}1
             {Ń}{{\'N}}1
}

\title{Postać kanoniczna Jordana}
\date{\today}

\begin{document}
\maketitle

Rozważamy przestrzeń $(\mathbb{X}^{n}, \mathbb{X}, +, \dot)$.% gdzie $\mathbb{X} = \mathbb{K}^{n}$

%%%%%%%%%%%%%%%%%%%%%%%%%%%%%%%%%%%%%%%%%%%%%%%%%%%%%%%%%%%%%%%%%%%%%%%%%%%%%%%%%%%%%%%%%
\section{Podprzestrzeń niezmiennicza względem macierzy $A$}
Podprzestrzeń $\mathbf{X} \subset \mathbb{X}^{n}$ nazywamy podprzestrzenią niezmienniczą względem danej macierzy $A$ (ogólniej względem przekształcenia) jeśli zachodzi warunek:
$$
    A\mathbf{X} \subseteq \mathbf{X},\, \text{tzn. } \forall_{x \in \mathbf{X}}{Ax \in \mathbf{X}}.
$$

%%%%%%%%%%%%%%%%%%%%%%%%%%%%%%%%%%%%%%%%%%%%%%%%%%%%%%%%%%%%%%%%%%%%%%%%%%%%%%%%%%%%%%%%%
\section{Najprostsza przestrzeń niezmiennicza}
Najprostszą podprzestrzenią niezmienniczą ze względu na macierz $A$ może być prosta.
Czyli podprzestrzeń, której baza składa się tylko z jednego wektora, który spełnia warunek:
$$
    A\mathbf{h} = \lambda\mathbf{w}
$$
Można zapisać go w postaci:
$$
    \left[ \lambda I -A \right]\mathbf{h} = \mathbf{0}
$$
Równanie to ma rozwiązanie $\mathbf{h} \neq \mathbf{0}$ wtedy i tylko wtedy gdy wyznacznik macierzy $\lambda I -A$ jest równy zeru:
$$
    |\lambda I -A| = 0,
$$
jest to tak zwane równanie charakterystyczne macierzy $A$.
Natomiast wielomian $W(\lambda) = |\lambda I -A|$ nazywamy wielomianem charakterystycznym macierzy $A$, a pierwiastki wielomianu charakterystycznego $\lambda_{i}$ nazywamy wartościami własnymi.
Ponadto zbiór wartości własnych macierzy $A$, oznaczany jako $\lambda(A) = \left\{ \lambda_{1}, \lambda_{2}, \cdots, \lambda_{k} \right\}$, nazywamy widmem macierzy.

%%%%%%%%%%%%%%%%%%%%%%%%%%%%%%%%%%%%%%%%%%%%%%%%%%%%%%%%%%%%%%%%%%%%%%%%%%%%%%%%%%%%%%%%%
\section{Diagonalizacja macierzy}
Można zauważyć, że jeśli macierz $A$ posiada $n$ wektorów własnych $\mathbf{w}_{1}, \mathbf{w}_{2}, \cdots, \mathbf{w}_{n}$ to można utworzyć z nich nieosobliwą macierz $P = \begin{bmatrix} \mathbf{w}_{1} & \mathbf{w}_{2} & \cdots & \mathbf{w}_{n} \end{bmatrix}$, wówczas:
$$
    \begin{array}{c}
        AP =
        A \begin{bmatrix} \mathbf{h}_{1} & \mathbf{h}_{2} & \cdots & \mathbf{h}_{n} \end{bmatrix}
        = \\
        \begin{bmatrix} \lambda_{1}\mathbf{h}_{1} & \lambda_{2}\mathbf{h}_{2} & \cdots & \lambda_{n}\mathbf{h}_{n} \end{bmatrix}
        =
        \begin{bmatrix} \mathbf{h}_{1} & \mathbf{h}_{2} & \cdots & \mathbf{h}_{n} \end{bmatrix}
        \begin{bmatrix}
            \lambda_{1} & 0             & 0 & \cdots & 0\\
            0           & \lambda_{2}   & 0 & \cdots & 0\\
            \vdots      & \vdots        & 0 & \cdots & \vdots\\
            0           & 0             & 0 & \cdots & \lambda_{n}
        \end{bmatrix}
        = P\Lambda.
    \end{array}
$$
A zatem:
$$
    AP = P\Lambda \implies \Lambda = P^{-1}AP.
$$
Oznacza to, że jeśli macierz posiada $n$ wektorów własnych, to jest podobna do macierzy diagonalnej.
Zatem istnieje baza, a dokładniej zestaw wektorów bazowych macierzy $A$, w której przekształcenie ma reprezentację w postaci macierzy diagonalnej.
Czyli najprostszą z możliwych reprezentacji.

%%%%%%%%%%%%%%%%%%%%%%%%%%%%%%%%%%%%%%%%%%%%%%%%%%%%%%%%%%%%%%%%%%%%%%%%%%%%%%%%%%%%%%%%%
\section{Co gdy nie istnieje $n$ wektorów własnych?}
W przypadku gdy macierz nie posiada $n$ wektorów własnych, nie można zbudować z nich macierzy przejścia między bazami i dokonać diagonalizacji.
Można jednak "dobrać" tzw. wektory główne.

Wektorem głównym, rzędu $m$ nazywamy wektor $\mathbf{h}^{m} \neq \mathbf{0}$ spełniający warunek:
$$
    (\lambda I - A)^{m}\mathbf{h}^{m} = \mathbf{0}, \, \text{oraz } (\lambda I - A)^{m-1}\mathbf{h}^{m} \neq \mathbf{0}.
$$
Na podstawie tej definicji można powiedzieć, że wektory własne są wektorami głównymi rzędu $1$.

Można wówczas zapisać, że zachodzi równość:
$$
    \begin{array}{c}
        (\lambda I - A)^{m} \mathbf{h}^{m} = \mathbf{0} = (\lambda I - A)^{m-1} \mathbf{h}^{m-1}\\
        (\lambda I - A) \mathbf{h}^{m} = \mathbf{h}^{m-1}\\
    \end{array}
$$
Na tej podstawie można  zapisać serię równości:
$$
    \begin{array}{l}
         A\mathbf{h}^{1} = \lambda\mathbf{h}^{1}\\
         A\mathbf{h}^{2} = \lambda\mathbf{h}^{2} + \mathbf{h}^{1}\\
         A\mathbf{h}^{3} = \lambda\mathbf{h}^{3} + \mathbf{h}^{2}\\
         \vdots \\
         A\mathbf{h}^{k} = \lambda\mathbf{h}^{k} + \mathbf{h}^{k-1}\\
    \end{array}
$$

Oznacza to, że możemy macierz $P$ utworzyć z wektorów własnych i wektorów głównych:
$$
    P
    =
    \begin{bmatrix}
        \mathbf{h}_{1}^{1} & \mathbf{h}_{1}^{2} & \cdots & \mathbf{h}_{1}^{m_{1}} &
        \mathbf{h}_{2}^{1} & \mathbf{h}_{2}^{2} & \cdots & \mathbf{h}_{2}^{m_{2}} &
        \cdots &
        \mathbf{h}_{k}^{m_{k}}
    \end{bmatrix}
$$

Wówczas można zapisać analogicznie jak w przypadku poprzednim:
$$
    \begin{array}{c}
         AP
         =
         A
        \begin{bmatrix}
            \mathbf{h}_{1}^{1} & \mathbf{h}_{1}^{2} & \cdots & \mathbf{h}_{1}^{m_{1}} &
            \mathbf{h}_{2}^{1} & \mathbf{h}_{2}^{2} & \cdots & \mathbf{h}_{2}^{m_{2}} &
            \cdots &
            \mathbf{h}_{k}^{m_{k}}
        \end{bmatrix} = \\
        \begin{bmatrix}
            A\mathbf{h}_{1}^{1} & A\mathbf{h}_{1}^{2} & \cdots & A\mathbf{h}_{1}^{m_{1}} &
            A\mathbf{h}_{2}^{1} & A\mathbf{h}_{2}^{2} & \cdots & A\mathbf{h}_{2}^{m_{2}} &
            \cdots &
            A\mathbf{h}_{k}^{m_{k}}
        \end{bmatrix}
        =\\
        \begin{bmatrix}
            \lambda_{1}\mathbf{h}_{1}^{1} \\(\mathbf{h}_{1}^{1} + \lambda_{1}\mathbf{h}_{1}^{2}) \\ \vdots \\ (\mathbf{h}_{1}^{m_{1}} + \lambda_{1}\mathbf{h}_{1}^{m_{1}-1}) \\
            \lambda_{2}\mathbf{h}_{2}^{1} \\ (\mathbf{h}_{2}^{1} + \lambda_{2}\mathbf{h}_{2}^{2}) \\ \vdots \\ (\mathbf{h}_{2}^{m_{2}} + \lambda_{2}\mathbf{h}_{2}^{m_{2}-1}) \\
            \vdots \\
            (\mathbf{h}_{k}^{m_{k}} + \lambda_{k}\mathbf{h}_{k}^{m_{k}-1})
        \end{bmatrix}^{T}
        =
        \\
        \begin{bmatrix}
            \mathbf{h}_{1}^{1} & \mathbf{h}_{1}^{2} & \cdots & \mathbf{h}_{1}^{m_{1}} &
            \mathbf{h}_{2}^{1} & \mathbf{h}_{2}^{2} & \cdots & \mathbf{h}_{2}^{m_{2}} &
            \cdots &
            \mathbf{h}_{k}^{m_{k}}
        \end{bmatrix} 
        \cdot\\
        \cdot
        \begin{bmatrix}
            \lambda_{1} &  1            & \cdots & 0        & 0         & \cdots & 0\\
            0           &  \lambda_{1}  & \cdots & 0        & 0         & \cdots & 0\\
            \vdots      &  \vdots       & \vdots & \vdots   & \vdots    & \vdots & \vdots\\
            0      &  0      & \cdots & 0   & 0    & \cdots & \lambda_{k}\\
        \end{bmatrix}=\\
        = PJ \implies J  = P^{-1}AP
    \end{array}
$$
Macierz $J$ nazywana jest macierzą Jordana, a jej struktura wygląda następująco:
$$
    J = 
    \left[
    \begin{array}{ccccc}
        J_{n_1}(\lambda_{1})    &   \mathbf{0}              &   \cdots  &   \mathbf{0}\\
        \mathbf{0}              &   J_{n_2}(\lambda_{2})    &   \cdots  &   \mathbf{0}\\
        \vdots                  &   \vdots                  &   \ddots  &   \vdots\\
        \mathbf{0}              &   \mathbf{0}              &   \cdots  &   J_{n_k}(\lambda_{k})
    \end{array}
    \right],
$$
gdzie macierz $J_{l}(\lambda)$, nazywana klatką Jordana związaną z wektorem własnym, jest macierzą wymiaru $n \times n$ o postaci;
$$
    J_{k}(\lambda)
    =
    \begin{bmatrix}
        \lambda & 1         & 0       &  \cdots & \cdots    & 0\\
        0       & \lambda   & 1       &  \cdots & \cdots    & 0\\
        \vdots  & \vdots    & \ddots  &  \ddots & \ddots    & \vdots\\
        0       & 0         & 0       &  \cdots & \lambda   & 1\\
        0       & 0         & 0       &  \cdots & 0         & \lambda
    \end{bmatrix}.
$$

%========================================================================================
\subsection{Własności macierzy Jordana}
\begin{itemize}
    \item Liczba $k$ klatek Jordana tworzących macierz $J$ jest równa liczbie liniowo niezależnych wektorów własnych macierzy $A$.
    \item Liczba klatek Jordana odpowiadających wartości własnej $\lambda$ jest równa liczbie odpowiadających jej liniowo niezależnych wektorów własnych.
    Suma stopni wszystkich tych klatek równa jest krotności wartości własnej $\lambda$ jako pierwiastka wielomianu charakterystycznego macierzy $A$.
    \item Liczba wektorów własnych związanych z wartością własną $\lambda$ jest równa wymiarowi jądra przekształcenia $(\lambda I - A)x$.
\end{itemize}

%========================================================================================
\subsection{Przykład}
Rozważmy macierz:
$$
    A 
    =
    \begin{bmatrix}
        2 & 0 &  0 & 0\\
        0 & 1 &  1 & 0\\
        0 & 0 &  1 & 0\\
        0 & 2 & -1 & 2
    \end{bmatrix}.
$$
Jej wielomian charakterystyczny wynosi:
$$
    W(\lambda) = |\lambda I  - A| = \lambda^{4} - 6\lambda^{3} + 13\lambda^{2} - 12\lambda + 4
    = (\lambda - 2)^{2}(\lambda - 1)^{2}.
$$
Ma ona zatem dwie podwójne wartości własne $\lambda_{1} = 1$ i $\lambda_{2}=2$.
Dla każdej z nich mamy:
\begin{itemize}
    \item Dla $\lambda_{1} = 1$ otrzymujemy $dim\{ Ker(1\cdot I-A)\} = 1$ a więc jeden wektor własny - konieczność wyznaczenia wektora głównego:
    \textit{Wektor własny}
    $$
        (1\cdot I-A)\mathbf{h}_{1}^{1} = \mathbf{0} \rightarrow
        \left[
        \begin{array}{cccc|c}
           -1 & 0 &  0 & 0 & 0\\
            0 & 0 & -1 & 0 & 0\\
            0 & 0 &  1 & 0 & 0\\
            0 &-2 &  1 &-1 & 0
        \end{array}
        \right]
        \rightarrow
        \mathbf{h}_{1}^{1}
        =
        \begin{bmatrix}
            0 \\ 1 \\ 0 \\ -2
        \end{bmatrix}
    $$
    \textit{Wektor główny}
    $$
        (1\cdot I-A)\mathbf{h}_{1}^{2} = \mathbf{h}_{1}^{1} \rightarrow
        \left[
        \begin{array}{cccc|c}
           -1 & 0 &  0 & 0 & 0\\
            0 & 0 & -1 & 0 & 1\\
            0 & 0 &  1 & 0 & 0\\
            0 &-2 &  1 &-1 &-2
        \end{array}
        \right]
        \rightarrow
        \mathbf{h}_{1}^{2}
        =
        a
        \begin{bmatrix}
            0 \\ 1 \\ 0 \\ -2
        \end{bmatrix}
        +
        \begin{bmatrix}
            0 \\ 0 \\ 1 \\ -11
        \end{bmatrix}
        \rightarrow
        \mathbf{h}_{1}^{2}
        =
        \begin{bmatrix}
            0 \\ 0 \\ 1 \\ -1
        \end{bmatrix}
    $$

    \item Dla $\lambda_{2} = 2$ otrzymujemy $dim\{ Ker(2\cdot I-A)\} = 2$ a więc dwa wektory własne:
    $$
        (2\cdot I-A)\mathbf{h} = \mathbf{0} \rightarrow
        \left[
        \begin{array}{cccc|c}
            0 & 0 &  0 & 0 & 0\\
            0 & 1 & -1 & 0 & 0\\
            0 & 0 &  1 & 0 & 0\\
            0 &-2 &  1 & 0 & 0
        \end{array}
        \right]
        \rightarrow
        \mathbf{h}
        =
        a
        \begin{bmatrix}
            1 \\ 0 \\ 0 \\ 0
        \end{bmatrix}
        +
        b
        \begin{bmatrix}
            0 \\ 0 \\ 0 \\ 1
        \end{bmatrix}
        \rightarrow
        \mathbf{h}_{2}
        =
        \begin{bmatrix}
            1 \\ 0 \\ 0 \\ 0
        \end{bmatrix},
        \mathbf{h}_{3}
        =
        \begin{bmatrix}
            0 \\ 0 \\ 0 \\ 1
        \end{bmatrix}
    $$
\end{itemize}
Mamy więc:
\begin{itemize}
    \item wektor własny $\mathbf{h}_{1}^{1}$ powiązany z klatką Jordana $J_{2}(1) = \begin{bmatrix} 1 & 1\\ 0 & 1 \end{bmatrix}$,
    \item wektor własny $\mathbf{h}_{2}^{1}$ powiązany z klatką Jordana $J_{1}(2) = \begin{bmatrix} 2 \end{bmatrix}$,
    \item wektor własny $\mathbf{h}_{2}^{2}$ powiązany z klatką Jordana $J_{1}(2) = \begin{bmatrix} 2 \end{bmatrix}$.
\end{itemize}

Zatem ostatecznie mamy:
\begin{itemize}
    \item 
    $
        P
        =
        \begin{bmatrix}
            0 & 0 & 1 & 0\\
            1 & 0 & 0 & 0\\
            0 & 1 & 0 & 0\\
           -2 &-1 & 0 & 1
        \end{bmatrix},
    $
    \item 
    $
        J
        =
        \begin{bmatrix}
            1 & 1 & 0 & 0\\
            0 & 1 & 0 & 0\\
            0 & 0 & 2 & 0\\
            0 & 0 & 0 & 2
        \end{bmatrix}.
    $
\end{itemize}

%%%%%%%%%%%%%%%%%%%%%%%%%%%%%%%%%%%%%%%%%%%%%%%%%%%%%%%%%%%%%%%%%%%%%%%%%%%%%%%%%%%%%%%%%
\section{Funkcje o argumentach macierzowych}
\subsection{Własności}
\begin{itemize}
    \item jeżeli $B = P^{-1}AP$, to zachodzi $f(B) = P^{-1}f(A)P$,
    \item jeżeli $A = diag(A_{1}, A_{2}, \cdots, A_{n})$, to zachodzi $f(A) = diag(f(A_{1}), f(A_{2}), \cdots, f(A_{n}))$,
\end{itemize}

%========================================================================================
\subsubsection{Wniosek}
Aby w prosty sposób wyznaczyć wartość funkcji $f(A)$ można znaleźć postać Jordana macierzy $A$, wyznaczyć $f(J)$ korzystając z drugiej własności i powrócić do oryginalnej bazy na podstawie własności pierwszej.

%========================================================================================
\subsubsection{Eksponenta macierzowa $e^{At}$}
Szczególnie ważną funkcją, ze względów praktycznych jest eksponenta macierzowa $e^{At}$.
Algorytm jej wyznaczania jest następujący:
\begin{enumerate}
    \item Wyznaczyć postać Jordana macierzy - $J$, oraz macierz modalną - $P$.
    \item Wyznaczyć eksponentę dla każdej z klatek Jordana $J_{n_{k}}(\lambda_{k})$:
    \begin{itemize}
        \item jeżeli wymiar klatki Jordana $n_{k} = 1$, to jest to przypadek skalarny: $e^{J_{n_{k}}t} = e^{\lambda_{k}t}$,
        \item jeżeli wymiar klatki Jordana $n_{k} > 1$, to postępujemy na podstawie następującego rozumowania:
            \begin{itemize}
                \item eksponentę rozwijamy w szereg Taylora: $e^{a} = \sum_{i=0}^{\infty}{\frac{a^n}{n!}}$,
                \item podstawiamy za $a$ macierz $J_{n_k}(\lambda_k)t$,
                \item rozkładamy $J_{n_k}(\lambda_k)t$ na sumę macierzy $I_{n_k}$ i $C_{n_k}(\lambda_k) = \begin{bmatrix} 0 & 1 & 0 & \cdots & 0\\ 0 & 0 & 1 & \cdots & 0\\ \vdots & \vdots & \vdots & \ddots & \vdots\\ 0 & 0 & 0 & \cdots & 0 \end{bmatrix}$,
                \item otrzymujemy wówczas $e^{J_{n_k}(\lambda_{k})t}  = e^{\lambda_{k}I_{n_k}t + C_{n_k}(\lambda_k)t} = e^{\lambda_{n_k}I_{n_k}t} \cdot e^{C_{n_k}(\lambda_k)t}$
                \item $e^{\lambda_{n_k}I_{n_k}t} =  \begin{bmatrix} e^{\lambda_{k}t} & 0 & 0 & \cdots & 0\\ 0 & e^{\lambda_{k}t} & 0 & \cdots & 0\\ \vdots & \vdots & \vdots & \ddots & \vdots\\ 0 & 0 & 0 & \cdots & e^{\lambda_{k}t} \end{bmatrix}$ - ponieważ jest to macierz diagonalna,
                \item zatem $e^{C_{n_k}(\lambda_{k})t}  = e^{C_{n_k}(\lambda_k)t} = I_{n_k} + C_{n_k}(\lambda_k)t + C_{n_k}(\lambda_k)^{2}t^2 + \cdots + C_{n_k}(\lambda_k)^{n_k -1}t^{n_k -1}$ - ponieważ macierz ściśle trójkątna w $n$-tej potędze się "zeruje",
                \item ostatecznie otrzymujemy $e^{J_{n_k}(\lambda_{k})t} = 
                \begin{bmatrix}
                    e^{\lambda_{k}t} & \frac{te^{\lambda_{k}t}}{1!} & \frac{t^2e^{\lambda_{k}t}}{2!} & \cdots & \frac{t^{n-1}e^{\lambda_{k}t}}{(n-1)!}\\
                    0                 & e^{\lambda_{k}t} & \frac{te^{\lambda_{k}t}}{1!} & \cdots & \frac{t^{n-2}e^{\lambda_{k}t}}{(n-2)!}\\
                    \vdots & \vdots & \vdots & \ddots & \vdots\\
                    0 & 0 & 0 & \cdots & e^{\lambda_{k}}
                \end{bmatrix}
                $
            \end{itemize}
        \item eksponenta macierzy Jordana ma więc postać:
        $
            e^{J}
            =
            \begin{bmatrix}
                e^{J_{1}(\lambda_{1})} & \mathbf{0} & \cdots & \mathbf{0}\\
                \mathbf{0} & e^{J_{2}(\lambda_{2})} & \cdots & \mathbf{0}\\
                \vdots & \vdots & \ddots & \vdots\\
                \mathbf{0} & \mathbf{0} & \cdots & e^{J_{k}(\lambda_{k})}
            \end{bmatrix}
        $
    \end{itemize}
    \item wyznaczamy $e^{At} = P e^{J} P^{-1}$.
\end{enumerate}

%========================================================================================
\subsubsection{Przykład wyznaczania $e^{At}$}
Rozważmy macierz:
$$
    A 
    =
    \begin{bmatrix}
        2 & 0 &  0 & 0\\
        0 & 1 &  1 & 0\\
        0 & 0 &  1 & 0\\
        0 & 2 & -1 & 2
    \end{bmatrix}.
$$

\begin{itemize}
    \item 
    $
        P = 
        \begin{bmatrix}
            1 & 0 &  0 & 0\\
            0 & 0 &  1 & \frac{1}{2}\\
            0 & 0 &  0 & 1\\
            0 & 1 & -1 &-2
        \end{bmatrix}
    $,
    $
        J = 
        \begin{bmatrix}
            2 & 0 &  0 & 0\\
            0 & 2 &  0 & 0\\
            0 & 0 &  1 & 1\\
            0 & 0 &  0 & 1
        \end{bmatrix}
    $,
    \item
    $
        \mathbf{h}_{2} \rightarrow e^{2t}
    $,
    $
        \mathbf{h}_{3} \rightarrow e^{2t}
    $,
    $
        \mathbf{h}_{1}^{2} \rightarrow
        e^{J_{2}(1)}
        =
        \left[\begin{array}{cc} {\mathrm{e}}^t & t\,{\mathrm{e}}^t\\ 0 & {\mathrm{e}}^t \end{array}\right]
    $
    \item 
    $
        e^{Jt}
        =
        \left[\begin{array}{cccc} {\mathrm{e}}^{2\,t} & 0 & 0 & 0\\ 0 & {\mathrm{e}}^{2\,t} & 0 & 0\\ 0 & 0 & {\mathrm{e}}^t & t\,{\mathrm{e}}^t\\ 0 & 0 & 0 & {\mathrm{e}}^t \end{array}\right]
    $
    \item
    $
        e^{At}
        =
        P e^{Jt} P^{-1}
        =
        \left[\begin{array}{cccc} {\mathrm{e}}^{2\,t} & 0 & 0 & 0\\ 0 & {\mathrm{e}}^t & t\,{\mathrm{e}}^t & 0\\ 0 & 0 & {\mathrm{e}}^t & 0\\ 0 & 2\,{\mathrm{e}}^{2\,t}-2\,{\mathrm{e}}^t & {\mathrm{e}}^{2\,t}-{\mathrm{e}}^t-2\,t\,{\mathrm{e}}^t & {\mathrm{e}}^{2\,t} \end{array}\right]
    $
\end{itemize}

\end{document}